\section*{Problem 1: Multi-Level Cache Hierarchy on an Intel Xeon Chip (21 points)}

\paragraph{Conventions used throughout.}
The statement's hint says that a miss penalty is an additional cost \emph{on top of} the
hit time, and that a miss is discovered only after the access to that level has been
attempted. The hierarchy is therefore modelled by the nested expression
\begin{equation}
\label{p1:eq:amat}
\mathrm{AMAT}=HT_{L1}+MR_{L1}\,P_{L1},\qquad
P_{L1}=HT_{L2}+MR_{L2}\,P_{L2},\qquad
P_{L2}=HT_{L3}+MR_{L3}\,P_{L3},
\end{equation}
with $HT_{L1}=1$, $HT_{L2}=10$, $HT_{L3}=50$ CC, and $P_{L3}$ the main-memory access
cost. The miss rates supplied in Questions C and D are \emph{local} miss rates, which is
exactly what \eqref{p1:eq:amat} consumes, so no conversion to global rates is needed.

For the memory access itself, the initial access time $M$ ``accounts for a single
transfer'', so with a block of $128$ B and a bus of $B$ bytes only the transfers after
the first are charged $10$ CC each:
\begin{equation}
\label{p1:eq:pl3}
P_{L3}=M+10\left(\frac{128\,\mathrm{B}}{B}-1\right).
\end{equation}
Finally, the clock is $2.0$ GHz, so $1\ \mathrm{CC}=1/(2.0\times10^{9}\,\mathrm{Hz})=0.5$ ns.

\subsection*{Question A (3 points)}

\begin{enumerate}
  \item[(1)] \textbf{False.} The L1 caches hold only $64$ KB each, so with L2 and L3
  removed every L1 miss would have to be served by main memory (hundreds of cycles)
  instead of by a $10$- or $50$-cycle cache, which raises the average memory access time
  substantially.

  \item[(2)] \textbf{True.} L1 must be fast enough to answer in a single clock cycle, so
  it is built from faster, lower-density SRAM cells with more ports and more decoding
  logic per byte; L3 is latency-tolerant and can use denser, cheaper cells, so a given
  capacity costs more in L1 than in L3.

  \item[(3)] \textbf{False.} Caches are filled on demand, one block at a time, and the
  program resides in main memory; only the blocks actually being referenced need to be
  present in L1. Indeed a program larger than $64$ KB could never fit in the L1
  instruction cache at all, yet still executes correctly.
\end{enumerate}

\subsection*{Question B (6 points)}

\paragraph{Block offset.}
Entries are \emph{Byte} addressable, so the offset field must select an individual byte
within a block. With a $128$-Byte block,
\[
  \#\text{offset bits}=\log_{2}128=7 .
\]
(The fact that a data word is $64$ bit $=8$ Bytes does \emph{not} shrink this field:
because addressing is per Byte and not per word, the low $3$ bits that select a byte
inside a word are still part of the address and belong to the offset.)

\paragraph{Index.}
The L2 cache holds $512\ \mathrm{KB}=512\times2^{10}=2^{19}$ Bytes, so the number of
blocks is
\[
  \#\text{blocks}=\frac{2^{19}\ \mathrm{B}}{128\ \mathrm{B}}=\frac{2^{19}}{2^{7}}=2^{12}=4096 .
\]
The cache is $4$-way set associative, i.e.\ $4$ blocks per set, so
\[
  \#\text{sets}=\frac{4096}{4}=1024=2^{10}
  \quad\Longrightarrow\quad
  \#\text{index bits}=\log_{2}1024=10 .
\]

\paragraph{Tag.}
The tag is whatever remains of the $64$-bit physical address:
\[
  \#\text{tag bits}=64-\#\text{index}-\#\text{offset}=64-10-7=47 .
\]

\paragraph{Answer.}
\begin{center}
\begin{tabular}{lccc}
\toprule
Field & Tag & Index & Block offset \\
\midrule
Width (bits) & $47$ & $10$ & $7$ \\
\bottomrule
\end{tabular}
\end{center}
Check: $47+10+7=64$ bits, as required.

\subsection*{Question C (10 points)}

Local miss rates are $MR_{L1}=5\%$, $MR_{L2}=40\%$, $MR_{L3}=20\%$, and the bus is
$16$ Bytes wide. A $128$-Byte block therefore needs $128/16=8$ transfers, of which
$8-1=7$ are \emph{additional}, so by \eqref{p1:eq:pl3}
\[
  P_{L3}=M+10\times 7=M+70 \quad\text{CC},
\]
where $M$ is the initial main memory access time in CCs. Substituting into
\eqref{p1:eq:amat}, the requirement $\mathrm{AMAT}\le 4$ CC reads
\[
  1+0.05\Bigl[\,10+0.40\bigl(50+0.20\,(M+70)\bigr)\Bigr]\;\le\;4 .
\]
Unwinding the nesting one level at a time:
\begin{align*}
  P_{L1}&\le\frac{4-1}{0.05}=60\ \text{CC},\\[2pt]
  P_{L2}&\le\frac{60-10}{0.40}=125\ \text{CC},\\[2pt]
  P_{L3}&\le\frac{125-50}{0.20}=375\ \text{CC},\\[2pt]
  M&\le 375-70=305\ \text{CC}.
\end{align*}

\paragraph{Conversion.}
At $2.0$ GHz one clock cycle is $0.5$ ns, hence
\[
  \boxed{\,M_{\max}=305\ \mathrm{CC}\times 0.5\ \frac{\mathrm{ns}}{\mathrm{CC}}=152.5\ \mathrm{ns}\,}
\]

\paragraph{Check.}
With $M=305$ CC: $P_{L3}=375$, $P_{L2}=50+0.2(375)=125$, $P_{L1}=10+0.4(125)=60$, and
$\mathrm{AMAT}=1+0.05(60)=4$ CC exactly, so the bound is attained and is tight.

\subsection*{Question D (2 points)}

The miss rates are unchanged and the initial main memory access time carried forward is
the maximum value found in Question C, $M=305$ CC ($=152.5$ ns). Only the bus width
changes, from $16$ to $32$ Bytes, so a $128$-Byte block now takes $128/32=4$ transfers,
of which $4-1=3$ are additional. By \eqref{p1:eq:pl3},
\[
  P_{L3}=305+10\times 3=335\ \text{CC}.
\]
Working back up through \eqref{p1:eq:amat}:
\begin{align*}
  P_{L2}&=50+0.20\times 335=117\ \text{CC},\\
  P_{L1}&=10+0.40\times 117=56.8\ \text{CC},\\
  \mathrm{AMAT}&=1+0.05\times 56.8=3.84\ \text{CC}.
\end{align*}
The average memory access time will be $\boxed{3.84}$ CCs (equivalently $1.92$ ns).
As expected, doubling the bus width removes $40$ CC of transfer time from the L3 miss
penalty and lowers the AMAT below the $4$ CC budget of Question C.
